\section{Bối cảnh và vấn đề}

\begin{frame}{Hệ thống HPC đang phục vụ workload ngày càng đa dạng}
\begin{columns}[T,onlytextwidth]
\begin{column}{0.53\textwidth}
\begin{itemize}
    \item Mô phỏng song song CPU và MPI.
    \item Huấn luyện, suy luận AI trên GPU.
    \item Xử lý dữ liệu quy mô lớn.
    \item Workflow nhiều giai đoạn.
    \item Yêu cầu về deadline, chi phí, ưu tiên và chất lượng dịch vụ.
\end{itemize}
\end{column}
\begin{column}{0.43\textwidth}
\begin{block}{Thực tế giao tiếp hiện nay}
Giao diện phổ biến giữa người dùng và scheduler vẫn là một batch script chứa yêu cầu tài nguyên cố định.
\end{block}
\begin{alertblock}{Hệ quả}
Tài nguyên không đồng nhất, nhưng ngữ nghĩa của nhu cầu người dùng lại bị nén vào một cấu hình duy nhất.
\end{alertblock}
\end{column}
\end{columns}
\end{frame}

\begin{frame}{Một resource request đang phải đại diện cho quá nhiều ý nghĩa}
\begin{columns}[T,onlytextwidth]
\begin{column}{0.46\textwidth}
\begin{block}{Ví dụ batch script}
\texttt{\#SBATCH --partition=gpu}\\
\texttt{\#SBATCH --gres=gpu:4}\\
\texttt{\#SBATCH --cpus-per-task=16}\\
\texttt{\#SBATCH --mem=64G}\\
\texttt{\#SBATCH --time=02:00:00}
\end{block}
\end{column}
\begin{column}{0.5\textwidth}
Một cấu hình duy nhất thường đồng thời được hiểu là:
\begin{enumerate}
    \item nhu cầu kỹ thuật tối thiểu;
    \item cấu hình mong muốn;
    \item giới hạn scheduler được phép chọn;
    \item lựa chọn tối ưu của job.
\end{enumerate}
\end{column}
\end{columns}
\vspace{0.3em}
Scheduler không biết 4 GPU là bắt buộc, bảo thủ, hay chỉ là một lựa chọn để đạt mục tiêu.
\end{frame}

\begin{frame}{Một request hợp lệ chưa chắc là một request hiệu quả}
\begin{block}{Intent thực tế}
“Tôi muốn job hoàn thành trước 18 giờ. Có thể giảm từ 4 xuống 2 GPU nếu vẫn kịp deadline và tiết kiệm quota.”
\end{block}
\begin{columns}[T,onlytextwidth]
\begin{column}{0.47\textwidth}
\textbf{Batch script chỉ biểu diễn}
\[
GPU=4,\qquad walltime=2\,h
\]
\end{column}
\begin{column}{0.47\textwidth}
\textbf{Intent còn chứa}
\begin{itemize}
    \item deadline;
    \item preference về chi phí;
    \item khoảng linh hoạt được cho phép;
    \item điều kiện chấp nhận cấu hình thay thế.
\end{itemize}
\end{column}
\end{columns}
\end{frame}

\begin{frame}{Khi intent không được biểu diễn, decision space bị thu hẹp}
\begin{columns}[T,onlytextwidth]
\begin{column}{0.47\textwidth}
\textbf{Đối với người dùng}
\begin{itemize}
    \item tăng thời gian chờ;
    \item tiêu thụ quota không cần thiết;
    \item bỏ lỡ deadline;
    \item phải thử--sai nhiều lần.
\end{itemize}
\end{column}
\begin{column}{0.47\textwidth}
\textbf{Đối với hệ thống}
\begin{itemize}
    \item giảm khả năng backfilling;
    \item giữ tài nguyên lớn cho job over-request;
    \item giảm mức độ đồng thời;
    \item tăng makespan, giảm utilization và throughput.
\end{itemize}
\end{column}
\end{columns}
\vspace{0.5em}
\begin{alertblock}{Vấn đề cốt lõi}
Scheduler không nhất thiết ra quyết định sai; scheduler chỉ được cung cấp một \textbf{decision space quá hẹp}.
\end{alertblock}
\end{frame}
